% Generated draft: requires mathematical and source review.
\documentclass[11pt]{article}
\usepackage[margin=25mm]{geometry}
\usepackage{fontspec}
\setmainfont{Latin Modern Roman}
\usepackage{amsmath,amssymb,mathtools,unicode-math}
\setmathfont{Latin Modern Math}
\usepackage{enumitem,needspace,xurl,hyperref}
\setlength{\parindent}{0pt}
\setlength{\parskip}{5pt}
\setlength{\emergencystretch}{3em}
\newcommand{\N}{\mathbb{N}}
\newcommand{\Z}{\mathbb{Z}}
\newcommand{\Q}{\mathbb{Q}}
\newcommand{\R}{\mathbb{R}}
\newcommand{\C}{\mathbb{C}}
\newcommand{\F}{\mathbb{F}}
\newcommand{\A}{\mathbb{A}}
\newcommand{\PP}{\mathbb{P}}
\newcommand{\E}{\mathbb{E}}
\newcommand{\eps}{\varepsilon}
\newcommand{\dd}{\,\mathrm{d}}
\begin{document}
\section*{UnsolvedMath ID: 10400066}
\textbf{Source problem code:} catalogue entry
\subsection{Definitions and mathematical statement}
The source asks for a geometric construction of ``these homology classes coming from Lie algebras'', then asks whether the formal-vector-field interpretation of graph cohomology can produce the Kontsevich integral on the topological knot side. The antecedent classes and their target homology group are not supplied.

For context, on a formal symplectic vector space \((V,\omega)\), a formal Hamiltonian \(H\) determines a formal vector field \(X_H\) by \(\iota_{X_H}\omega=\dd H\); Hamiltonians modulo constants have Lie bracket given by the Poisson bracket. Continuous Lie-algebra cohomology of formal vector fields provides the Gelfand--Fuchs framework. The graph-complex correspondence may involve stable dimension, relative cohomology and restrictions on Hamiltonians, none of which is specified in the excerpt.

The requested construction should use knots together with the appropriate algebras of formal vector fields, and yield the Kontsevich integral \(Z(K)\) in the intended completed diagram algebra, with equality in every degree after matching conventions. This is stronger than merely applying a Lie-algebra weight system to an already constructed \(Z(K)\). The source also asks for geometric realization of the indicated Lie-algebra-derived homology classes, but the classes, degrees and realization space must be recovered before that part becomes precise.

The passage attributes to Kontsevich an algebraic explanation/generalization of Rozansky--Witten weight systems through this cohomology and asks for a topological counterpart. No exact theorem, Lie algebra, graph differential, or knot-space construction is supplied. The summary calls the topological construction open/triage while acknowledging the algebraic interpretation; these are source attributions pending current primary-literature review.
\subsection{Short English statement}
Realize the indicated Lie-algebra homology classes geometrically, and use formal Hamiltonian vector fields with knot geometry to recover the Kontsevich integral. The relevant classes, precise Lie/graph complex and topological construction are missing from the excerpt.
\subsection{Sources}
\begin{itemize}
\item Local registry: \url{data/unsolvedmath/problems.json}, permanent UnsolvedMath ID \(10400066\), source code catalogue entry; supplied statement, background, and research summary, updated 2026-08-06.
\item Ohtsuki, \emph{Problems on invariants of knots and 3-manifolds} (2002), Problem 3.16, attributed to R. Bott, PDF page 65: \url{https://msp.org/gtm/2002/04/gtm-2002-04-024s.pdf}. Supplied PDF-font-structure extraction accessed 2026-07-29; provenance distinguishes 231 numbered items.
\item Kontsevich, Rozansky--Witten and Gelfand--Fuchs are named by the passage without full original-paper citations or an exact specified correspondence.
\end{itemize}
\end{document}
